\documentclass[11pt,letterpaper]{article}
\usepackage[utf-8]{inputenc}
\usepackage[margin=1in]{geometry}
\usepackage{hyperref}
\usepackage{url}
\usepackage{graphicx}
\usepackage{booktabs}
\usepackage{amsmath}
\usepackage{amssymb}

\title{Why Activation-Level Speculative Decoding Can't Work:\\A Measurement-Driven Negative Result}

\author{Webmind Research}

\date{April 16, 2026}

\begin{document}

\maketitle

\begin{abstract}
Activation-level speculative decoding hypothesizes that predicting the next hidden state at a shard boundary can hide network latency in distributed transformer inference. We measure this hypothesis on a real distributed WebGPU LLM system (Synapse) deployed across heterogeneous mobile devices. Using GPT-2 117M with a 2-shard topology and 96 measurement samples from coherent decode runs, we find that a linear-extrapolation predictor achieves mean cosine similarity of $0.38 \pm 0.12$ against ground-truth activations. Zero samples reach an acceptance threshold of $0.7$; the maximum observed was $0.607$. This is not a tuning problem: consecutive autoregressive activations do not move locally linearly, because each new token injects a fresh low-rank update whose direction is precisely the computation downstream decoders are built to perform, not extract from prior steps. We show that this limitation is architectural, present even with longer prediction history, and that the only viable routes forward are learned predictors or token-level draft-verify schemes. This is a tier-b reproducible-methodology negative result, with full code paths and shadow-mode instrumentation disclosed.
\end{abstract}

\section{Introduction}

Distributed transformer inference across network-connected devices faces a fundamental bottleneck: the ratio of network round-trip latency ($\sim 50$\,ms) to on-GPU compute time for a single token forward pass ($\sim 1$\,ms). On mobile-dominated topologies, the wire, not the compute, limits throughput.

The CPU branch-prediction analogy is compelling. In pipelined processors, branch prediction speculatively executes down the predicted path, rolling back on misprediction—the cost of rollback is masked by the speculative work. Here, \emph{activation prediction} at a shard boundary would let a downstream node compute speculatively on the predicted hidden state while the true activation is in flight. Correct prediction hides latency; incorrect prediction costs only one recompute—no worse than the baseline.

This paper tests the core hypothesis: are consecutive autoregressive transformer activations predictable from their local history?

\section{The Seductive Hypothesis}

We hypothesized that a linear-extrapolation predictor,
\[\text{next} \approx \text{curr} + (\text{curr} - \text{prev}),\]
should produce useful signal if autoregressive decode moves smoothly through activation space. The predictor is trivially cheap (one vector add per step), requires only two observations, and carries an explicit confidence estimate via the cosine of consecutive deltas.

\section{Bootstrap Deadlock: A Pattern to Avoid}

Before measurement, we hit a gating problem we created ourselves.

The controller was designed to prove itself before activation:
\[\text{enabled} = \text{true} \Leftarrow \text{verified\_steps} \geq \text{warmup\_threshold}\]
\[\Leftarrow \text{verify()} \text{ calls}\]
\[\Leftarrow \text{pending speculations exist}\]
\[\Leftarrow \text{speculation kicked off}\]
\[\Leftarrow \text{enabled} = \text{true}\]

This is circular gatekeeping. On a live deploy with 59 cached decode steps, we logged zero speculation events. The feature was dead code from the start.

The fix was to allow \texttt{\_speculateNext()} to run in shadow mode during warmup, with acceptance still gated on \texttt{enabled}. But shadows still produced zero predictions.

A second gate was hiding: the predictor silently dropped predictions with confidence $< 0.9$. In early decode, the cosine of consecutive deltas hovers around $0.5$–$0.85$; most predictions never reached the \texttt{verify()} logging surface.

\textbf{Same bootstrap-deadlock pattern, one layer deeper.}

The general lesson: any self-enabling optimization with a \emph{prove yourself before I turn you on} threshold risks nested gates. The fix is cheap shadow-mode instrumentation—run predict and verify without gating, discard the output, keep only the diagnostic statistic. This adds one cosine computation per step and produces unbiased measurement without corrupting inference.

\section{Measurement}

\subsection{Setup}

Test rig: GPT-2 117M, 2-shard split (layers 0–5 / 6–11), homogeneous Qualcomm mobile topology (Adreno GPU via Android Chrome).

Protocol:
\begin{enumerate}
\item Run 40-token autoregressive decode at temperature $T=0.9$.
\item At each step, compute $\text{cosine}(\text{curr\_activation}, \text{predicted\_activation})$ in shadow mode.
\item Log the cosine without gating on prediction confidence.
\item Repeat with two different prompts; 96 samples total.
\end{enumerate}

\subsection{Results}

\begin{table}[h]
\centering
\begin{tabular}{ll}
\toprule
$n$ & 96 \\
Mean cosine & 0.380 \\
Stdev & 0.124 \\
Min & -0.158 \\
p25 & 0.309 \\
p50 (median) & 0.388 \\
p75 & 0.469 \\
p95 & 0.564 \\
Max & 0.607 \\
\bottomrule
\end{tabular}
\end{table}

\begin{table}[h]
\centering
\begin{tabular}{lrr}
\toprule
Threshold & Count / 96 & Percentage \\
\midrule
$\geq 0.99$ & 0 & 0\% \\
$\geq 0.90$ & 0 & 0\% \\
$\geq 0.80$ & 0 & 0\% \\
$\geq 0.70$ & 0 & 0\% \\
$\geq 0.60$ & 1 & 1.0\% \\
\bottomrule
\end{tabular}
\end{table}

The distribution is tight, unimodal, and sits entirely below 0.7. The hypothesis's implied acceptance threshold of 0.995 is roughly four standard deviations above the best-case observed cosine. Speculation would fire \emph{never} in production.

\subsection{Cosine Does Not Improve with Longer History}

A natural follow-up: maybe 2-step linear extrapolation is too short-sighted. Perhaps with longer history (exponential moving average or sliding-window regression), cosine climbs.

The data contradicts this. Bucketed by sequence position:

\begin{table}[h]
\centering
\begin{tabular}{lrrrr}
\toprule
Seq. Position & $n$ & Mean Cosine & Stdev & Max \\
\midrule
[5, 10) & 3 & 0.232 & 0.152 & 0.343 \\
[10, 15) & 5 & 0.424 & 0.076 & 0.497 \\
[15, 20) & 5 & 0.385 & 0.115 & 0.585 \\
[20, 25) & 5 & 0.321 & 0.081 & 0.415 \\
[25, 30) & 5 & 0.358 & 0.027 & 0.390 \\
[30, 35) & 5 & 0.352 & 0.049 & 0.439 \\
[35, 40) & 5 & 0.333 & 0.037 & 0.382 \\
[40, 45) & 5 & 0.295 & 0.052 & 0.342 \\
\bottomrule
\end{tabular}
\end{table}

Cosine peaks at sequence position 10–15 (shortly after prefill), then \emph{declines monotonically} to $\sim 0.30$ by position 40. Activations become \emph{less} predictable as generation progresses, not more.

\section{Why This Had to Happen}

At decode step $t$, the GPT-2 activation at the shard-1 boundary is shaped by one thing: the token sampled at step $t-1$. That token is a one-hot index into the 50,257-element vocabulary, drawn from the model's output distribution over the previous hidden state.

The sampled tokens at steps $t-1$ and $t-2$ point in essentially unrelated directions (GPT-2's embedding table is $\sim 64$-dimensional-effective in 768-dimensional space). After six transformer layers of non-linearity, the activation-space motion between consecutive steps is dominated by whichever one-hot direction the previous sampler picked, modulated by attention to the KV cache.

Linear extrapolation of that trajectory is equivalent to predicting the next sampled token's embedding from the current token's embedding—which is precisely what the remainder of the transformer is built to compute, and what you cannot obtain for free by arithmetic on the previous two hidden states.

The limitation is not tuning. It is architectural.

\section{What Would Work}

\begin{enumerate}
\item \textbf{Learned activation predictor.} Distill a small (e.g., 2-layer, 256-dim) MLP trained on captured activation sequences from real decode traces. A non-linear predictor can exploit attention-structure regularities that linear extrapolation misses. Cosine will still be bounded by the information content of the next token, but room for improvement exists.

\item \textbf{Token-level draft-verify.} The canonical approach: a small fast model proposes $K$ tokens ahead; the big model verifies them in a single forward pass. MEDUSA \cite{medusa}, SpecDec \cite{specdec}, and Lookahead Decoding \cite{lookahead} all follow this pattern. It sidesteps the shard-boundary bottleneck entirely because it does not rely on activation smoothness.

\item \textbf{Multi-token prediction heads.} Train the big model itself to predict $K$ tokens from one hidden state, so a single forward pass produces multiple accepted tokens.
\end{enumerate}

\section{Generalized Lesson: Nested Gates in Self-Enabling Optimizers}

The dual-deadlock pattern is a code-review pattern worth remembering. If you build an adaptive-precision selector, dynamic batcher, JIT-threshold tuner, or any self-promoting optimization with a "warmup required" gate, audit the code path that is supposed to produce the warmup data. Ensure it cannot be gated on the very thing it is supposed to enable.

Shadow-mode (zero-cost instrumentation running alongside the feature, with output discarded but statistics logged) is the general fix.

\section{Reproducibility}

Code paths in Synapse (https://github.com/tejasphatak/Synapse):
\begin{itemize}
\item \texttt{synapse-src/node/speculative.js} — shadow-mode predict + verify; \texttt{lastShadowCosine} field.
\item \texttt{synapse-src/node/node.js} — \texttt{speculation\_cosine\_sample} log emission.
\item \texttt{synapse-src/node/predictor.js} — linear-extrapolation predictor under test.
\item \texttt{synapse-src/coordinator/index.js} — log-store and \texttt{/api/logs} query surface.
\end{itemize}

Query live measurements directly:
\begin{verbatim}
curl -s 'https://coord/api/logs?event=speculation_cosine_sample' \
  | jq -r '.logs[].data.cosine'
\end{verbatim}

Results are deployment-dependent on hardware and prompt. The tight distribution centered in $[0.3, 0.5]$ reproduced across every configuration we tested.

\section{Acknowledgements}

This finding was produced on the Synapse distributed-inference prototype. Measurement was possible only because the cluster runs on browser WebGPU across heterogeneous mobile and desktop hardware in realistic network conditions—not a synthetic benchmark. The shadow-mode instrumentation design was informed by prior work on self-enabling optimizers in CPU microarchitecture. Contributions welcome at https://github.com/tejasphatak/Synapse.

\begin{thebibliography}{99}

\bibitem{medusa}
  Cai, T., Li, Z., Geng, Z., et al. (2024).
  \newblock ``Medusa: Simple LLM Inference Acceleration Framework with Multiple Decoding Heads.''
  \newblock \textit{arXiv preprint arXiv:2401.10891}.

\bibitem{specdec}
  Leviathan, Y., Kalman, Y., \& Matias, Y. (2023).
  \newblock ``Fast Transformer Decoding: One Write-Head is All You Need.''
  \newblock \textit{arXiv preprint arXiv:1911.02727}.

\bibitem{lookahead}
  Fu, Y., Prabhavalkar, R., Qian, K., et al. (2024).
  \newblock ``Lookahead Decoding.''
  \newblock \textit{arXiv preprint arXiv:2402.02057}.

\end{thebibliography}

\end{document}
